\section*{Sets and Indices}

\begin{itemize}
    \item $I$ : set of activities, here
    \[
    I=\{A,B,C,D,E,F,G\}.
    \]
    \item $P \subseteq I \times I$ : set of precedence arcs $(i,j)$, where activity $i$ must finish before activity $j$ can start. The arcs are read from \texttt{table\_1.csv}.
\end{itemize}

\section*{Parameters}

All numerical values are read from \texttt{general\_parameters.csv} and \texttt{table\_1.csv}.

\begin{itemize}
    \item $d_i$ : nominal duration of activity $i$ in days (code source: \texttt{dur[a]} from fields \texttt{activity\_A\_duration}, \ldots, \texttt{activity\_G\_duration}).
    \item $c^{\text{work}}$ : project work cost per day (field \texttt{work\_cost\_per\_day}).
    \item $c^{\text{mach}}$ : machine rental cost per day (field \texttt{machine\_rental\_cost\_per\_day}).
    \item $r^{E}$ : duration reduction for activity $E$ if overtime is used (field \texttt{activity\_E\_overtime\_reduction}).
    \item $c^{E,\text{ot}}$ : fixed overtime premium for accelerating activity $E$ (field \texttt{activity\_E\_overtime\_cost}).
    \item $c^{E,\text{rev}}$ : fixed follow-up review cost triggered when activity $E$ is accelerated (field \texttt{activity\_E\_followup\_review\_cost}).
\end{itemize}

For compactness, define the realized duration
\[
\tilde d_i =
\begin{cases}
d_E - r^E y, & i=E,\\
d_i, & i\neq E,
\end{cases}
\]
where $y$ is the overtime decision defined below.

\section*{Decision Variables}

\begin{itemize}
    \item $S_i$ (code: \texttt{S[a]}) : start time of activity $i$, with $S_i \ge 0$.
    \item $T$ (code: \texttt{T}) : project completion surrogate, with $T \ge 0$.
    \item $y$ (code: \texttt{use\_ot}) : binary variable equal to 1 if overtime is used for activity $E$, and 0 otherwise.
\end{itemize}

\section*{Objective}

The implemented model minimizes
\[
\min \quad
c^{\text{work}} T
\;+\;
c^{\text{mach}} \bigl(S_B + d_B - S_A\bigr)
\;+\;
\bigl(c^{E,\text{ot}} + c^{E,\text{rev}}\bigr)y.
\]

This exactly matches the code logic:
\begin{itemize}
    \item work cost is charged on $T$,
    \item machine rental cost is charged over the interval from the start of $A$ to the finish of $B$, namely $S_B + d_B - S_A$,
    \item accelerating activity $E$ is a gate decision: if $y=1$, both the overtime premium and the follow-up review cost are incurred; if $y=0$, neither is incurred.
\end{itemize}

\section*{Constraints}

\paragraph{Precedence constraints.}
For every $(i,j)\in P$,
\[
S_j \;\ge\; S_i + \tilde d_i.
\]

Equivalently, for this instance:
\[
\begin{aligned}
S_G &\ge S_A + d_A,\\
S_D &\ge S_A + d_A,\\
S_F &\ge S_E + (d_E-r^E y),\\
S_F &\ge S_G + d_G,\\
S_C &\ge S_D + d_D,\\
S_C &\ge S_F + d_F,\\
S_B &\ge S_F + d_F.
\end{aligned}
\]

\paragraph{Completion-time constraints.}
The code defines $T$ only through the finishing times of activities $C$ and $B$:
\[
T \ge S_C + d_C,
\qquad
T \ge S_B + d_B.
\]

\paragraph{Variable domains.}
\[
S_i \ge 0 \qquad \forall i\in I,
\]
\[
T \ge 0,
\]
\[
y \in \{0,1\}.
\]

\section*{Notes on Implementation}

\begin{itemize}
    \item The formulation is a mixed-integer linear program solved by the implementation with \texttt{GUROBI\_CMD}.
    \item The realized duration of activity $E$ is modeled exactly as in the code:
    \[
    \tilde d_E = d_E - r^E y.
    \]
    Thus overtime shortens $E$ by 3 days when selected.
    \item The revision requirement is implemented as a binary-triggered fixed cost, not as a proportional surcharge. In the objective, choosing overtime for $E$ adds
    \[
    c^{E,\text{ot}} + c^{E,\text{rev}},
    \]
    where the review cost corresponds to the required follow-up safety review.
    \item The model does not use the parameter \texttt{activity\_E\_followup\_review\_trigger} explicitly; the trigger is encoded directly by charging the review cost whenever \texttt{use\_ot}=1.
    \item The machine rental term uses $d_B$ (nominal duration of $B$) rather than any project-wide makespan expression, and $T$ is constrained only by activities $B$ and $C$. This is not simplified here; it is written to match the current implementation exactly.
\end{itemize}
